\section{Background}

\subsection{Coupling transforms}

A \emph{coupling transform} $\bm{\phi}$ \citep{Dinh:2015:nice} maps an input $ \bfx $ to an output $ \bfy $ in the following way: 
\begin{enumerate}[topsep=0pt, itemsep=1pt]
    \item Split the input $ \bfx $ into two parts, $ \bfx = \squarebr{\bfx_{1:d-1}, \bfx_{d:D}} $.
    \item Compute parameters $ \bftheta = \text{NN}(\bfx_{1:d-1}) $, where NN is an arbitrary neural network. 
    \item Compute $ y_i = g_{\bftheta_i}(x_i) $ for $i=d,\dots,D$ in parallel, where $g_{\bftheta_i}\!$ is an invertible function parameterized by $ \bftheta_i $.
    \item Set $ \bfy_{1:d-1} = \bfx_{1:d-1} $, and return $ \bfy = \squarebr{\bfy_{1:d-1}, \bfy_{d:D}} $.
\end{enumerate}
The Jacobian matrix of a coupling transform is lower triangular, since $ \bfy_{d:D} $ is given by transforming $ \bfx_{d:D} $ elementwise as a function of $ \bfx_{1:d-1} $, and $ \bfy_{1:d-1} $ is equal to $ \bfx_{1:d-1} $. 
Thus, the Jacobian determinant of the coupling transform $\bm{\phi}$ is given by 
$\det\roundbr{\frac{\partial\bm{\phi}}{\partial\bfx}}   = \prod_{i=d}^{D} \frac{\partial g_{\bftheta_i}}{\partial x_i}$, the product of the diagonal elements of the Jacobian.

Coupling transforms solve two important problems for normalizing flows: they have a tractable Jacobian determinant, and they can be inverted exactly in a single pass. 
The inverse of a coupling transform can be easily computed by running steps 1--4 above, this time inputting $ \bfy $, and using $g_{\bftheta_i}^{-1} $ to compute $ \bfx_{d:D} $ in step 3. 
Multiple coupling layers can also be composed in a natural way to construct a normalizing flow with increased flexibility. A coupling transform can also be viewed as a special case of an autoregressive transform where we perform two splits of the input data instead of $ D $, as noted by \citet{Papamakarios:2017:maf}. In this way, advances in flows based on coupling transforms can be applied to autoregressive flows, and vice versa.

\subsection{Invertible elementwise transformations}
\label{sec:invertible-elementwise-transformations} 
\paragraph{Affine/additive} Typically, the function $ g_{\bftheta_i} \!$ takes the form of an \textit{additive} \citep{Dinh:2015:nice} or \textit{affine} \citep{Dinh:2017:rnvp} transformation for computational ease. The affine transformation is given by:
\begin{equation}
    g_{\bftheta_i}(x_i) = \alpha_i x_i + \beta_i,
    \quad\text{where}\quad
    \bftheta_i = \curlybr{\alpha_{i}, \beta_{i}}, 
\end{equation}
and $ \alpha_i $ is usually constrained to be positive. 
The additive transformation corresponds to the special case $\alpha_i\!=\!1$.
Both the affine and additive transformations are easy to invert, but they lack flexibility.
Recalling that the base distribution of a flow is typically simple, flow-based models may struggle to model multi-modal or discontinuous densities using just affine or additive transformations, since they may find it difficult to compress and expand the density in a suitably nonlinear fashion (for an illustration, see \cref{appendix-2d-datasets-affine}). 
We aim to choose a more flexible $g_{\bftheta_i} $, that is still differentiable
and easy to invert.

\paragraph{Polynomial splines} Recently, \citet{Muller:2018:nis} proposed a powerful generalization of the above affine transformations, based on monotonic piecewise polynomials. 
The idea is to restrict the input domain of $g_{\bftheta_i}\! $ to the interval $ [0, 1] $, partition the input domain into $ K $ bins, and define $ g_{\bftheta_i}\! $ to be a simple polynomial segment within each bin. 
\citet{Muller:2018:nis} restrict themselves to monotonically-increasing linear and quadratic polynomial segments, whose coefficients are parameterized by $\bftheta_i$. 
Moreover, the polynomial segments are restricted to match at the bin boundaries so that  $g_{\bftheta_i} \!$ is continuous. Functions of this form, which interpolate between data using piecewise polynomials, are known as \emph{polynomial splines}.

\paragraph{Cubic splines} 
In a previous iteration of this work
\ifanonymize\citep{anonymous:2019:cubicsplineflows}\else\citep{Durkan:2019:cubicsplineflows}\fi,
we explored the \emph{cubic-spline flow}, a natural extension to the framework of \citet{Muller:2018:nis}. 
We proposed to implement $g_{\bftheta_i}\!$ as a \emph{monotonic cubic spline} \cite{steffen1990simple}, where $g_{\bftheta_i}\!$ is defined to be a monotonically-increasing cubic polynomial in each bin.
By composing coupling layers featuring elementwise monotonic cubic-spline transforms with invertible linear transformations, we found flows of this type to be much more flexible than the standard coupling-layer models in the style of RealNVP \citep{Dinh:2017:rnvp}, achieving similar results to autoregressive models on a suite of density-estimation tasks. 

Like \citet{Muller:2018:nis}, our spline transform and its inverse were defined only on the interval $ [0, 1] $. 
To ensure that the input is always between $0$ and $1$, we placed a sigmoid transformation before each coupling layer, and a logit transformation after each coupling layer. 
These transformations allow the spline transform to be composed with linear layers, which have an unconstrained domain.
However, the limitations of 32-bit floating point precision mean that in practice the sigmoid saturates for inputs outside the approximate range of $ [-13, 13] $, which results in numerical difficulties.
In addition, computing the inverse of the transform requires inverting a cubic polynomial, which is prone to numerical instability if not carefully treated \citep{blinn2007solve}. 
In \cref{sec:rq-splines} we propose a modified method based on rational-quadratic splines which overcomes these difficulties.

\subsection{Invertible linear transformations}
\label{sec:invertible-linear-transformations}
To ensure all input variables can interact with each other, it is common to randomly permute the dimensions of intermediate layers in a normalizing flow. 
Permutation is an invertible linear transformation, with absolute determinant equal to $ 1 $.
\citet{oliva2018tan} generalized permutations to a more general class of linear transformations, and \citet{Kingma:2018:glow} demonstrated improvements on a range of image tasks.
In particular, a linear transformation with matrix $\bfW$ is parameterized in terms of its \textit{LU-decomposition} $\bfW = \bfP \bfL \bfU$, where $ \bfP $ is a fixed permutation matrix, $ \bfL $ is lower triangular with ones on the diagonal, and $ \bfU $ is upper triangular. 
By restricting the diagonal elements of $\bfU$ to be positive, $ \bfW $ is guaranteed to be invertible.

By making use of the LU-decomposition, both the determinant and the inverse of the linear transformation can be computed efficiently.
First, the determinant of $ \bfW $ can be calculated in $\bigo{D}$ time as the product of the diagonal elements of $\bfU$. 
Second, inverting the linear transformation can be done by solving two triangular systems, one for $\bfU$ and one for $\bfL$, each of which costs $ \mathcal{O}(D^{2} M) $ time where $ M $ is the batch size. 
Alternatively, we can pay a one-time cost of $ \mathcal{O}(D^{3}) $ to explicitly compute $ \bfW^{-1} $, which can then be cached for re-use.
